\section{Political-Corruption Classifier Diagnostics}
\label{app:political-corruption-classifier}

\begin{table}
\caption{Full validation comparison of political-corruption classification models.}
\label{tab:pc-classifier-comparison-appendix}
\resizebox{\textwidth}{!}{%
\begin{tabular}{lllrrrrrrrrr}
\toprule
Run & Model & Labels & N train & Thr. & Acc. & PC Prec. & PC Rec. & PC F1 & Macro F1 & Wtd. F1 & Pred. pos. \\
\midrule
comparison & E5-large & Silver 1+2 & 3982 & 0.400 & 0.863 & 0.731 & 0.809 & 0.768 & 0.835 & 0.865 & 0.311 \\
comparison & E5-large & Silver 2 & 2995 & 0.400 & 0.849 & 0.704 & 0.794 & 0.747 & 0.819 & 0.851 & 0.317 \\
comparison & E5-large & Silver 1 & 987 & 0.450 & 0.843 & 0.701 & 0.766 & 0.732 & 0.810 & 0.845 & 0.307 \\
comparison & E5-large & Human CV & 502 & 0.500 & 0.817 & 0.634 & 0.823 & 0.716 & 0.790 & 0.823 & 0.365 \\
comparison & TF-IDF char. n-grams & Human CV & 502 & 0.500 & 0.777 & 0.601 & 0.610 & 0.606 & 0.725 & 0.777 & 0.285 \\
\bottomrule
\end{tabular}
}
\par\smallskip\footnotesize{Note. PC = political corruption. Precision, recall, and PC F1 refer to the political-corruption class.}
\end{table}


\begin{table}
\caption{Decision-threshold sweep for the selected political-corruption classifier.}
\label{tab:pc-classifier-threshold-sweep}
\resizebox{\textwidth}{!}{%
\begin{tabular}{rrrrrrrr}
\toprule
Thr. & Acc. & PC Prec. & PC Rec. & PC F1 & Macro F1 & Wtd. F1 & Pred. pos. \\
\midrule
0.150 & 0.612 & 0.419 & 0.993 & 0.589 & 0.610 & 0.620 & 0.665 \\
0.200 & 0.713 & 0.495 & 0.979 & 0.657 & 0.705 & 0.726 & 0.556 \\
0.250 & 0.747 & 0.528 & 0.943 & 0.677 & 0.734 & 0.760 & 0.502 \\
0.300 & 0.787 & 0.577 & 0.908 & 0.705 & 0.769 & 0.797 & 0.442 \\
0.350 & 0.829 & 0.649 & 0.851 & 0.736 & 0.805 & 0.835 & 0.369 \\
0.400 & 0.863 & 0.731 & 0.809 & 0.768 & 0.835 & 0.865 & 0.311 \\
0.450 & 0.855 & 0.762 & 0.702 & 0.731 & 0.816 & 0.853 & 0.259 \\
0.500 & 0.849 & 0.793 & 0.624 & 0.698 & 0.799 & 0.843 & 0.221 \\
0.550 & 0.833 & 0.820 & 0.518 & 0.635 & 0.763 & 0.819 & 0.177 \\
0.600 & 0.829 & 0.867 & 0.461 & 0.602 & 0.746 & 0.810 & 0.149 \\
0.700 & 0.797 & 0.933 & 0.298 & 0.452 & 0.663 & 0.756 & 0.090 \\
\bottomrule
\end{tabular}
}
\par\smallskip\footnotesize{Note. PC = political corruption. Precision, recall, and PC F1 refer to the political-corruption class.}
\end{table}


\begin{table}
\caption{Country-level validation performance of the selected political-corruption classifier.}
\label{tab:pc-classifier-country-validation}
\resizebox{\textwidth}{!}{%
\begin{tabular}{lrrrrrrrrr}
\toprule
Country & N & PC support & Acc. & PC Prec. & PC Rec. & PC F1 & Macro F1 & Wtd. F1 & Pred. pos. \\
\midrule
Sweden & 51 & 6 & 1.000 & 1.000 & 1.000 & 1.000 & 1.000 & 1.000 & 0.118 \\
France & 49 & 13 & 0.939 & 0.812 & 1.000 & 0.897 & 0.927 & 0.941 & 0.327 \\
Hungary & 51 & 21 & 0.882 & 0.826 & 0.905 & 0.864 & 0.880 & 0.883 & 0.451 \\
Serbia & 52 & 12 & 0.904 & 0.818 & 0.750 & 0.783 & 0.860 & 0.902 & 0.212 \\
Bulgaria & 48 & 18 & 0.833 & 0.778 & 0.778 & 0.778 & 0.822 & 0.833 & 0.375 \\
Ukraine & 48 & 23 & 0.750 & 0.739 & 0.739 & 0.739 & 0.750 & 0.750 & 0.479 \\
Italy & 51 & 12 & 0.863 & 0.692 & 0.750 & 0.720 & 0.815 & 0.865 & 0.255 \\
UK & 101 & 26 & 0.802 & 0.583 & 0.808 & 0.677 & 0.767 & 0.811 & 0.356 \\
Netherlands & 51 & 10 & 0.843 & 0.600 & 0.600 & 0.600 & 0.751 & 0.843 & 0.196 \\
\bottomrule
\end{tabular}
}
\par\smallskip\footnotesize{Note. PC = political corruption. Precision, recall, and PC F1 refer to the political-corruption class.}
\end{table}


\section{Political-Corruption Attention Descriptives}
\label{app:political-corruption-attention}

\input{output/tables/attention/table_attention_country_year_matrix}

\input{output/tables/attention/table_attention_peak_months}

\begin{figure}
\centering
\includegraphics[width=\textwidth]{output/figures/attention/political_corruption_absolute_volume_month_stacked.png}
\caption{Monthly absolute volume of articles classified as political corruption by country. The stacked areas report the number of articles classified as primarily about political corruption in each country-month. Unlike the relative-attention measure in the main text, these counts are not normalized by total news coverage and therefore reflect differences in both attention and corpus size.}
\label{fig:pc-absolute-volume-appendix}
\end{figure}

Figure~\ref{fig:pc-absolute-volume-appendix} reports absolute monthly volume as
a descriptive complement to the relative-attention measure used in the main
analysis.

\clearpage
\section{Exploratory Topic Model}
\label{app:exploratory-topic-model}

\input{output/tables/topic_models/topic_model_manuscript_values}

The exploratory topic model was estimated on a country-year-stratified sample
of \TopicModelSampleN{} articles from the final source-filtered
political-corruption corpus ($N=\TopicModelClassifierN{}$; classifier threshold
$=\TopicModelClassifierThreshold{}$). BERTopic identified
\TopicModelFineTopicN{} non-outlier fine-grained clusters. Of the sampled
articles, \TopicModelInlierN{} (\TopicModelInlierShare\%) were assigned to
these clusters and \TopicModelOutlierN{} (\TopicModelOutlierShare\%) were
treated as outliers.

GPT-5.1 generated descriptive labels for the fine-grained clusters from their
keywords and representative documents. The researchers then defined
\TopicModelConfiguredGroupN{} broad interpretive groupings, and GPT-5.1
assigned each fine-grained cluster to one grouping. These groups summarize how
political-corruption coverage is organized in this corpus; they are not
validated corruption-type measures and are not used as inferential outcomes.
Table~\ref{tab:topic_higher_order_summary} reports their weighted distribution.

\input{output/tables/topic_models/table_topic_higher_order_summary}

The complete fine-grained topic inventory, prompts, representative examples,
raw model responses, assignments, interactive visualizations, and checksummed
run manifests are retained in the Research Drive topic-model archive. The full
inventory is archived rather than inserted into the manuscript by default.
